\chapter{Structured memory, locality and timescales}

A model reads a long DNA record. It must preserve an exact recent motif,
track a slowly changing compositional signal, and avoid mixing one record
with the next. Should all three jobs belong to the same vector? A single
state can be useful, but naming it memory does not specify which
information it retains, which queries it can answer, or when it must reset.

Chapter 19 defined one selectively updated DOGMA reference. This chapter
separates memory representation, retention and update clocks. We construct
a small differentiable bank and a block-rate summary, prove their behavior,
and expose failures that a final-state equality test can miss. DOGMA remains
a proposed non-Transformer DNA-native research direction. The code below is
a reference operator, not a trained genomic model or a biological mechanism.

\section{Begin with the query a memory must answer}

Three requests illustrate different contracts. Returning the last four
canonical bases requires an ordered local buffer. Estimating a weighted
recent feature average can use a compressed state. Retrieving the feature
associated with a particular earlier location needs some notion of
addressing and retention of that association. A long decay constant does
not by itself provide an addressable record.

Fix an input feature sequence $u_t\in\mathbb R^D$. The features might come
from Chapter 19's DNA-pair table, but here we take them as explicit inputs
so memory behavior can be tested separately from the encoder. The reference
state will contain an exponential bank $H\in\mathbb R^{K\times D}$, a coarse
summary $s\in\mathbb R^D$, a pending sum $p\in\mathbb R^D$, and an integer
phase $q$. A recent raw-token buffer is a contrasting design, not an
unimplemented feature that the code secretly provides.

\Plate{01-layout}{Memory layout expresses different promises. Ordered raw
token slots support exact recent recall; the numerical bank supports weighted
summaries. The implemented reference contains the bank, coarse state, pending
sum and phase, not the illustrative token buffer.}{fig:layout}

\Term{Locality} must also be named precisely. Adjacent genomic coordinates,
neighboring positions in an encoded stream, adjacent tensor addresses, and
recent update events are different neighborhoods. A tokenizer that combines
several bases changes the relation between token distance and base-pair
distance. Packed records can even make two adjacent tensor slots belong to
different chromosomes or organisms. Memory ownership cannot be inferred
from storage adjacency.

\section{Derive a timescale from a declared clock}

For one channel, consider the relaxation equation
$dh/d\xi=(u-h)/\tau$, where $\xi$ is a declared clock coordinate and
$u$ is constant over an interval of length $\Delta\geq0$.
Here $\tau>0$ has the same units as $\Delta$. Solving the equation gives
\begin{equation}
 h_{\rm new}=a h_{\rm old}+(1-a)u,\qquad
 a=e^{-\Delta/\tau}.
 \label{eq:relax}
\end{equation}
The clock could count valid tokens, genomic distance, or elapsed time.
The equation alone does not choose among them. For a valid-token clock,
we use $\Delta=1$ on each valid feature and do nothing for padding.
An unobserved stretch of genome is not automatically padding: if distance
matters, it may need a nonzero $\Delta$ and a separately specified input policy.

The \Term{e-fold time} is $\tau$: a perturbation of the incoming state is
multiplied by $e^{-1}$ after this much elapsed clock distance under a fixed
coefficient. The half-life is $\tau\log 2$. For constant per-event retention
$a\in(0,1)$, $\tau=-1/\log a$. These are sensitivity timescales, not measured
biological half-lives or guaranteed task-relevant context lengths.

\Listing{retention}
Computing the write weight as $-\operatorname{expm1}(-\Delta/\tau)$ avoids
subtracting two nearly equal numbers when the interval is very short.
In exact arithmetic it equals $1-a$. In floating point the separately
computed coefficients may not sum to exactly one; tests use stated numerical
tolerances. We do not promise bitwise equality across arbitrary kernels.

\begin{center}\input{\Generated/timescales.tex}\end{center}
Retention $7/8$ forgets an existing perturbation more slowly than $1/2$,
but its immediate response to a fresh input is only $1/8$ rather than $1/2$.
Longer retention trades immediate write amplitude for persistence in this
normalized update. Simply increasing retention is not a free improvement.

\section{Unroll the state and inspect what it remembers}

For a fixed $a$ and initial $h_0$,
\begin{equation}
 h_t=a^t h_0+\sum_{j=1}^{t}(1-a)a^{t-j}u_j.
 \label{eq:kernel}
\end{equation}
The proof is substitution: multiply the previous expression by $a$ and add
$(1-a)u_t$. The input at lag $\ell=t-j$ has weight
$w_\ell=(1-a)a^\ell$, with $\sum_{\ell=0}^{\infty}w_\ell=1$.
For a constant input and zero initial state, $h_t=(1-a^t)u$.
For one unit impulse followed by zeros, the response is exactly the kernel.

\Plate{02-kernel}{Two generated impulse kernels use the same unit input and
zero initial state. Lag zero is the state immediately after that input.
The slow bank writes less initially and retains a longer tail; neither curve
is a genomic accuracy measurement.}{fig:kernel}

With $K$ timescales, update row $k$ of the bank as
$H_{k,t}=a_kH_{k,t-1}+(1-a_k)u_t$. Each row summarizes the same feature
history with a different kernel. A readout such as
$y_t=\sum_k w_kH_{k,t}$ can combine those summaries. It cannot recover
information that every row has discarded. If the readout weights are
unconstrained, the readout need not be a convex average even though every
bank row is one.

\begin{proposition}[A finite bank is a bounded summary]
If every input coordinate and initial bank coordinate lies in $[-M,M]$,
then every bank coordinate remains in that interval in exact arithmetic.
For two runs with the same inputs, their difference in row $k$ after $t$
unit-clock events is $a_k^t$ times its initial difference.
\end{proposition}
\begin{proof}
Each update is a convex combination, proving the range bound by induction.
Subtract the recurrences of the two runs; the input terms cancel.
Iterating $\delta H_{k,t}=a_k\delta H_{k,t-1}$ gives the second claim.
\end{proof}

The second statement concerns perturbations of the initial state, not
arbitrary parameter gradients or separability of every pair of histories.
Chapter 19 already showed how a state-dependent gate can invalidate this
simple contraction argument. A bank of fixed timescales is deliberately
chosen here so the memory kernel is inspectable.

\section{Slow updates are not the same as slow retention}

A second mechanism updates a coarse state only after a block of $B$ valid
features. Within the current block retain a sum $p$ and phase $q$.
For each valid input,
\begin{equation}
 p\leftarrow p+u_t,\qquad q\leftarrow q+1.
\end{equation}
If $q=B$, commit
\begin{equation}
 s\leftarrow\rho s+(1-\rho)\frac{p}{B},\qquad
 p\leftarrow0,\quad q\leftarrow0,
 \quad 0\leq\rho\leq1.
 \label{eq:commit}
\end{equation}
Otherwise $s$ is held exactly. Padding changes none of these fields.
The block boundary is determined by valid-event count, not by how a loader
cuts a record into transport chunks.

After $m$ complete blocks, an initial coarse-state perturbation is
$\rho^m$ times its original value. At block-aligned positions,
an equivalent per-token envelope is $a_{\rm eff}=\rho^{1/B}$.
Between commits, however, the state is constant. Thus a per-token exponential
with retention $a_{\rm eff}$ is not the same operator: it responds earlier,
weights positions inside a block differently, and has different gradient paths.

\Plate{03-clock}{A chunk ending after position 3 carries phase 3 and the
pending sum. The fourth valid position completes the block and triggers the
coarse update. The transport boundary does not reset or flush the model's
clock. A read before the commit cannot use a future block summary.}{fig:clock}

For $B=4$, $\rho=1/2$, and input $(1,0,0,0,0,0,0,0)$, the first block
mean is $1/4$ and the coarse state becomes $1/8$. It remains there until
the next commit, then becomes $1/16$. The two exponential rows respond on
every valid position instead.
\begin{center}\input{\Generated/trace.tex}\end{center}

What happens to a partial block at the end of a record? Our reference retains
it without committing. A caller may inspect it as part of the complete state,
or discard it when the record is finished. An explicit finalization operation
could average by the actual count, but that would be a different output
contract. Flushing every transport chunk is not valid finalization.

\Needspace{12\baselineskip}
\section{Implement the complete carried state}

\Listing{Memory}
The container's fields name the sufficient state for this recurrence.
The frozen dataclass prevents rebinding its attributes; it does not make
PyTorch tensors deeply immutable. The implementation performs no in-place
updates, and callers must not mutate tensors retained for later continuation
or gradient computation.

Before executing, the validator checks feature shape $[T,D]$, positive
timescales $[K]$, shared dtype/device, a Boolean mask, the allowed phase, and
the shapes of all carried tensors. An empty pending block must have zero sum.
It checks even masked values for finiteness so NaNs are not silently hidden.
This is a correctness-oriented reference, not a device-efficient validation
kernel.
\Listing{multiscale}

The helper \texttt{validate} is included in full in the companion source,
along with a separately expanded convolution oracle. The build extracts
whole functions for printed listings. No line-number fragment is maintained
independently of the executable file.

The returned history includes the initial state and one state per input slot.
It is convenient for teaching and testing, but storing this history costs
$O(TKD)$ values. A streaming inference implementation that retains only the
terminal state can have bounded carried storage; that is not a claim that
training through the reference uses constant memory.

\section{Prove chunk parity, then test the missing-field failure}

Each valid event maps $(H,s,p,q)$ deterministically to its successor using
fixed configuration $(\tau,B,\rho)$. A masked slot is the identity map.
Applying the same sequence of maps with a temporary pause therefore gives
the same final state. This proves chunk parity in exact arithmetic,
provided all four fields cross the boundary and the configuration is unchanged.

The proof is short because the state is complete. If only $H$ and $s$
are saved after three valid positions of a four-position block, the next
chunk cannot know whether its first token should trigger a commit.
Dropping $p,q$ changes the model even when a casual observer sees plausible
bank values. The tests deliberately perform this broken continuation and
check that both the coarse state and phase differ.

\Plate{05-ownership}{Record A's complete carried state follows only record A.
Record B begins from its own reset state. Dashed forbidden edges mark
cross-record reuse. In a serving system the record and model identities are
additional ownership metadata, not implicit tensor properties.}{fig:ownership}

Forward parity is only half of the training contract. A saved tensor that
has been detached from the autograd graph may retain the right values but
lose dependence on earlier inputs. Tests compare gradients into both
features and timescales across a split. For independent records, resetting
is correct; for a continuing record, resetting is a semantic change.
Persistent storage must also bind model revision, dtype, timescales, block
size, readout convention and any upstream pair-encoder state.

\section{Coarse memory loses order: an exact counterexample}

Take $B=2$ and zero initial coarse state. The blocks $(1,-1)$ and $(-1,1)$
both have mean zero. They leave the same coarse state for every $\rho$.
Yet a fast bank with $a=1/2$ ends at $-1/4$ and $+1/4$, respectively.
Therefore the coarse state cannot answer a question about which sign came
last, while the fast bank can distinguish these particular examples.

\Plate{04-aliasing}{A block mean erases within-block order. The two histories
meet at the same coarse value but reach different fast values. Adding a
fast state repairs this counterexample; it does not prove lossless recovery
of every possible history.}{fig:aliasing}

This is \Term{aliasing} under the selected summary map, not numerical noise.
No more accurate floating-point arithmetic can make the two exact block
means different. A learned feature encoder might preserve relevant
distinctions before pooling, but that is a hypothesis requiring a task and
evidence. Declaring a summary global memory does not settle its adequacy.

A separate counting argument applies to any finite-precision implementation.
If all carried state and metadata together have at most $C$ bits, there are
at most $2^C$ representable states. There are $4^n=2^{2n}$ canonical DNA
strings of length $n$. Whenever $2n>C$, two distinct strings must share a
state. This proves a limit on lossless full-history encoding, not failure
on every task. Many tasks need only a small sufficient statistic.
Ideal exact real numbers do not satisfy the finite-bit premise.

\section{Timescale gradients and useful information}

For one update with $a=e^{-\Delta/\tau}$ and fixed incoming $h,u$,
\begin{equation}
 \frac{\partial a}{\partial\tau}
 =a\frac{\Delta}{\tau^2},\qquad
 \frac{\partial h_{\rm new}}{\partial\tau}
 =a\frac{\Delta}{\tau^2}(h-u).
 \label{eq:taugrad}
\end{equation}
Increasing $\tau$ moves the update toward the old state. If old state and
input coincide, this local derivative vanishes. Across time, the total
derivative must include earlier states' dependence on $\tau$; treating
each step's incoming state as a constant misses those paths.

For a block commit, each of its $B$ features has derivative
$(1-\rho)/B$ into the new coarse state before later attenuation.
A subsequent $m$ commits multiply this direct contribution by $\rho^m$.
Before its own block commits, the feature has no effect on $s$, although
it is present in $p$ and in the fast bank. Reading a completed block summary
back into earlier positions would leak future information.

Long retention preserves a sensitivity but does not select which information
deserves preservation. A model may retain nuisance composition while losing
a rare functional pattern. Conversely, short memory can be enough when
the target depends only on local sequence. Report performance by declared
dependency distance and perturbation type instead of turning $\tau$ into
an unsupported effective-context-length headline.

\section{Research comparisons without architectural substitution}

HiPPO formulates recurrent memory through online polynomial approximation
of a history under a specified measure \cite{hippo}. Its approximation
objective gives a more structured meaning to compression than an arbitrary
collection of decay constants. Our independent exponential rows are not
a HiPPO construction, and inherit no optimality theorem from that work.

Titans combines a learned test-time memory with other memory roles
\cite{titans}. In its long-term module, an associative objective drives
updates to memory parameters, with a momentum-like state and forgetting.
Our recurrence updates activation state with fixed parameters during a
record; it is not test-time learning. The shared lesson is to specify
what is stored, how it is written and what a read returns, not to treat
all mechanisms called memory as equivalent.

The distinction matters for DNA-domain systems. Molecular binding, cellular
regulation, neural activation state, optimizer state and a persistent
database have different units and causal mechanisms. A software timescale
measured in valid tokens is not a cellular lifetime. Biological labels
should not stand in for mechanism diagrams when the object is a tensor.

\section{Build an evaluation contract that could reject the design}

Compare at least a local-window baseline, a single-timescale recurrence,
the multiscale bank, and the bank plus block summary. Hold the feature
encoder and readout policy fixed when isolating memory effects. Match or
report parameter count, carried-state bytes, arithmetic, training budget
and actual hardware latency separately. A broader trained baseline may be
needed to test usefulness; a controlled ablation tests a narrower causal claim.

Choose the split by the scientific question. Related sequences and overlapping
windows should not cross a purportedly independent train/test boundary.
State whether genomic coordinates, species, homologous groups or individuals
are held out. Report more than an average score: probe rare patterns,
long gaps, changed composition, ambiguous bases, record boundaries and
reverse-complement transformations with the appropriate label semantics.
Chapter 21 will own the strand-symmetry contract; a directional recurrence
does not become strand invariant by being DNA-addressed.

For this reference, persistent numerical state contains $(K+2)D$ scalars
plus phase. At 32 bits per scalar, $K=8,D=64$ needs
$4(8+2)64=2560$ bytes for those tensors, excluding allocator overhead,
parameters, history, metadata and any raw-token buffer.
Each valid event costs $O(KD)$ arithmetic for the bank and $O(D)$ for the
pending sum. Coarse commits add $O(D)$ every $B$ valid events.
These counts are not measured throughput.

Operational tests should interleave records, change chunk boundaries, insert
padding, serialize and restore complete state, and deliberately supply a
mismatched model revision. Our local reference tests arithmetic, causality,
masking and in-memory ownership; it does not implement a versioned serving
store. Industrial claims require that additional system and its failure tests.

\section{Exercises with worked answers}
\begin{enumerate}
\item Convert $a=1/2$ into a unit-clock e-fold time and half-life.
\textbf{Answer:} $\tau=1/\log2$; half-life $1$ valid event.
\item Starting from zero, apply two unit inputs with $a=3/4$.
\textbf{Answer:} $h_1=1/4$, $h_2=7/16=1-(3/4)^2$.
\item What is the weight of an input at lag three for $a=1/2$?
\textbf{Answer:} $(1-a)a^3=1/16$, not $1/8$.
\item Why is two-clock composition exact only under a stated input policy?
\textbf{Answer:} With the same held input,
$a(\Delta_1)a(\Delta_2)=a(\Delta_1+\Delta_2)$ and the write weights combine.
Different inputs produce different weighted contributions.
\item With $B=4,\rho=1/2$, what follows block $(1,0,0,0)$ from reset?
\textbf{Answer:} $s=1/8$, $p=0$, $q=0$; the fast bank has its own values.
\item A chunk ends with $q=3$. May the next chunk start with $q=0$?
\textbf{Answer:} Not for continuation. The next valid event must commit the
old pending sum plus the new feature.
\item Show that equal block means need not imply equal fast states.
\textbf{Answer:} $(1,-1)$ and $(-1,1)$ have mean zero; at $a=1/2$ their
terminal fast values are $-1/4$ and $+1/4$.
\item Differentiate the update with respect to $\tau$ at fixed $h,u$.
\textbf{Answer:} Equation~\ref{eq:taugrad}; its sign follows $h-u$.
Total sequence derivatives also propagate through previous states.
\item Can a $C$-bit state losslessly represent all length-$n$ DNA strings
when $2n>C$? \textbf{Answer:} No, by the pigeonhole principle.
This does not rule out exact task-specific sufficient statistics.
\item Does a frozen dataclass prevent in-place tensor mutation?
\textbf{Answer:} No. It prevents attribute rebinding, not mutation of a
referenced tensor. Ownership rules and non-mutating updates remain necessary.
\item Compute carried tensor bytes for $K=4,D=32$ at float32.
\textbf{Answer:} $4(4+2)32=768$ bytes, before phase and other overheads.
\item Propose a test of useful long-range memory.
\textbf{Rubric:} Fix the task and label alignment; separate related training
and test sequences; use matched baselines and encoder/readout controls;
vary dependency distance and nuisance context; report uncertainty, state
cost, runtime and negative cases. Synthetic kernel plots alone earn no
credit as evidence of genomic capability.
\end{enumerate}

\section{Vocabulary and the next boundary}
\Term{Update clock}: the events or distance increments that advance state.
\Term{Retention}: attenuation of previous state in a declared recurrence.
\Term{Carried state}: all fields required to resume the same computation.
\Term{Information capacity}: the distinctions representable under the
chosen precision and storage budget, not merely the length of input read.
The next chapter asks how these directional memory states should transform
when the same DNA is represented on the opposite strand.
